\SourceLeaf{68}
Confectum est igitur, quatenus per series fieri potuit, negotium propositum.

Cognata est disquisitio, quam Ill. \textsc{Kl\"ugel} (de functione potentiali

\[
(1-2x\cos\varphi+xx)^{-m}
\]

\noindent in Commentat. Societatis huius Tom. XII) exhibuit. Forma, quam ille contemplatus
est, abludit quidem ab illa, quae ducente \textsc{Eulero} in hac dissertatiuncula in censum
vocata fuit, verum difficile non foret, alteram ex altera deducere. Ne igitur acta
omnino egisse videar, moneam, omnem disquisitionem eum in finem a me institutam,
ut serierum infinitarum tractationem per terminos generales nulla difficultate laborare
et facilius saepe, ac artificia singularia, ad eruendas leges generales perducere,
exemplo aliquo ostenderem. Non igitur in principali dissertatiunculae parte usus sum
via, qua VV. Ill. \textsc{Euler} et \textsc{Kl\"ugel} incessere; in ceteris omnia ad expressiones generales
revocare et demonstratione ubicunque munire conatus.

\textls[140]{Corollarium.}

11. Occurrere in Analysi altiore series, sive finitas, sive infinitas, quarum termini
non unam eandemque legem primo obtutu sequi videntur, imprimis eiusmodi, in
quibus termini indicis paris aliter ac imparis constructi sunt, notissimum est. Solent
tunc singulari lege hi, singulari illi circumscribi. Praeter prolixitatem autem eiusmodi
definitionum, inest illis aliquid contra naturam ac vim Analyseos. In eo enim ipso
huius consistit dignitas, ut valores unius quantitatis, si eiusmodi fuerint, diversissimos
sub una expressione generali comprehendat, et ex hac expressione cunctos derivare
doceat. Sic, si in problemate aliquo geometrico solvendo plures una quantitate
solutionem praestent, hae omnes in aequatione, qua datarum cum quaesita relationem
circumscribimus, involutae reperiuntur. Simile quid ab expressionibus, quibus serierum
termini generales sistuntur, postulari potest. Si fuerit igitur inter terminos pares et
impares aliqua diversitas, ita comparata sit oportet, si perfecta fuerit, termini generalis
formula, ut alios valores sumto pro termini indice numero indefinito pari, alios pro
impari praebeat. Quod qua ratione praestari possit, iam in prooemio dissertatiunculae
exposui. Est nimirum $(-1)^m$ existente $m$ pari $=+1$, existente impari $=-1$; estque adeo
